% ALT TEXT: Four cells are shown inside one dashed relational-system boundary. A is the origin, with arrows A to B, A to C, and B to D; C and D are co-present at the same depth, giving A greater than B greater than the tied pair C equals D. A note says that one system cell carries the separately established observer/reference role but that this theorem does not identify that cell with B. Below, the strict restriction A before B before D maps to minimum before middle before maximum, forcing R of B to equal middle. Another note records 4m equals 2 over pi and warns that ancestry D is not the geometric filled share D.
\begin{figure}[t]
\centering
\begin{tikzpicture}[
  node distance=8mm and 16mm,
  cell/.style={regular polygon,regular polygon sides=6,draw=blue!65!black,
    fill=blue!5,minimum size=11mm,inner sep=1pt,font=\small},
  interior/.style={regular polygon,regular polygon sides=6,
    draw=orange!80!black,fill=orange!16,very thick,
    minimum size=13mm,inner sep=1pt,font=\small,align=center},
  rank/.style={draw,rounded corners,minimum width=16mm,inner sep=3pt},
  arrow/.style={-{Latex[length=2mm]},thick}
]
\node[cell] (A) {\(A\)};
\node[interior, below left=10mm and 15mm of A] (B)
  {\(B\)\\[-1pt]\scriptsize interior};
\node[cell, below right=10mm and 15mm of A] (C) {\(C\)};
\node[cell, below right=10mm and 15mm of B] (D) {\(D\)};
\draw[arrow] (A)--(B);
\draw[arrow] (A)--(C);
\draw[arrow] (B)--(D);
\draw[densely dashed] (C)--node[below,font=\scriptsize]{co-present}(D);
\node[draw=blue!45,densely dashed,rounded corners,
  fit=(A)(B)(C)(D),inner sep=5mm] (system) {};
\node[font=\scriptsize,fill=white,inner sep=1pt] at (system.north east)
  {four-cell relational system};
\node[above=2mm of A] {\(A>B>\{C=D\}\)};
\node[rank, below=18mm of B] (amin) {\(A\mapsto\min\)};
\node[rank, right=of amin,draw=orange!80!black,very thick]
  (bmid) {\(B\mapsto\operatorname{mid}\)};
\node[rank, right=of bmid] (dmax) {\(D\mapsto\max\)};
\draw[arrow] (amin)--node[above]{\(\prec\)}(bmid);
\draw[arrow] (bmid)--node[above]{\(\prec\)}(dmax);
\node[draw,rounded corners,below=6mm of bmid,inner sep=4pt]
  {\(\boxed{\mathcal R(B)=\operatorname{mid}}\)};
\node[align=left,font=\scriptsize,anchor=west] at (5.4,1.72)
  {One designated system cell carries the separately\\
   established observer/reference role.  This placement\\
   theorem does not identify that cell with \(B\).\\[3pt]
   \(4m=\dfrac2\pi\) at the completed four-cell location.\\
   Here ancestry \(D\) is a cell label.\\
   It is not the geometric \(D=\dfrac{\pi}{3\sqrt2}\).};
\end{tikzpicture}
\caption{Observer-rooted ancestry and middle placement.  The observer is the
separately established role of one cell within the relational system, not an
external eye or frame; this diagram does not assign that role to \(B\).  The
invariant lineage arrows are \(A\to B\), \(A\to C\), and \(B\to D\), with
\(C\) and \(D\) co-present.  Betweenness on the restricted chain
\(A\prec B\prec D\) forces the ancestry-interior element \(B\) to the middle
rank under the lawful observer readout.  The ancestry label \(D\) and the
geometric filled-share symbol \(D\) are distinguished by type; the displayed
four-cell composite does not identify their numerical values.}
\label{fig:observer-betweenness}
\end{figure}
